% from ega3-1-fr.tex
\oldpage[0\textsubscript{III}]{76}
On peut en outre supposer que ce système projectif est strict. Pour
cela, on observe qu'en vertu de l'isomorphisme
(\hyperref[0.13.7.1.1-fr]{13.7.1.1}), $A_k$ est isomorphe à
$A_{k+1}/\mathrm F^{k+1}(A_{k+1})$~; on peut donc prendre, dans la
construction de $X_{k+1}^\bullet$, la résolution injective de
$A_{k+1}/\mathrm F^{k+1}(A_{k+1})$ égale à $X_k^\bullet$, et il résulte
de (M, V, 2.3) que la construction du morphisme de complexes
$X_{k+1}^\bullet\to X_k^\bullet$ peut se faire de sorte que ce morphisme
fournisse par passage aux quotients un isomorphisme
\[
X_{k+1}^\bullet/\mathrm F^{k+1}(X_{k+1}^\bullet)
\xrightarrow{\sim}X_k^\bullet
\]
respectant les filtrations, ce qui est la condition de
(\hyperref[0.13.5.1-fr]{13.5.1}).
\end{env}

\begin{env}[13.7.3]
\label{0.13.7.3-fr}
Par construction, le système projectif $\mathrm T(\mathbf X^\bullet)$
des complexes $\mathrm T(X_k^\bullet)$ satisfait aux hypothèses de
(\hyperref[0.13.5.3-fr]{13.5.3}). Les résultats de
(\hyperref[0.13.5.4-fr]{13.5.4}),
(\hyperref[0.13.5.5-fr]{13.5.5}) et
(\hyperref[0.13.5.6-fr]{13.5.6}) sont donc applicables aux suites
spectrales $\mathrm E(\mathrm T(X_k^\bullet))=\mathrm E(A_k)$~; nous
écrirons $\mathrm E_r^{pq}(\mathbf A)$ au lieu de
$\mathrm E_r^{pq}(\mathrm T(\mathbf X^\bullet))$ pour
$1\leq r\leq+\infty$ (cf. (\hyperref[0.13.5.7-fr]{13.5.7}) pour
$r=+\infty$) et de même pour les notations analogues. On notera en
particulier que l'on a
\begin{equation}
\label{0.13.7.3.1-fr}
\tag{13.7.3.1}
\mathrm E_1^{pq}(\mathbf A)=
\mathrm R^{p+q}\mathrm T(\operatorname{gr}^p(\mathbf A))
\end{equation}
en vertu de (\hyperref[0.13.6.4.2-fr]{13.6.4.2}) et du fait que le
système $(\operatorname{gr}^p(A_k))$ est essentiellement constant.

Ces résultats et (\hyperref[0.13.4.3-fr]{13.4.3}) donnent la proposition
suivante, démontrée d'abord par Shih Weishu par une méthode différente
(inédite) :
\end{env}

\begin{proposition}[13.7.4]
\label{0.13.7.4-fr}
\emph{(Shih). --- Soit $n$ un entier. Les deux conditions suivantes sont
équivalentes :}
\begin{enumerate}
\item[\emph{a)}] \emph{Pour tout couple $(p,q)$ tel que $p+q=n$, le
système projectif $(\mathrm E_r^{pq}(\mathbf A))_{r\geq2}$ est
essentiellement constant.}
\item[\emph{b)}] \emph{Le système projectif
$\mathrm R^n\mathrm T(\mathbf A)=
(\mathrm R^n\mathrm T(A_k))_{k\in\mathbb Z}$ vérifie (ML).}
\end{enumerate}
\emph{En outre, lorsque ces conditions sont remplies, on a un
isomorphisme canonique}
\begin{equation}
\label{0.13.7.4.1-fr}
\tag{13.7.4.1}
\operatorname{gr}^p(\mathrm R^n\mathrm T(\mathbf A))
\xrightarrow{\sim}\mathrm E_\infty^{p,n-p}(\mathbf A)
\qquad\emph{pour tout $p\in\mathbb Z$.}
\end{equation}
\end{proposition}
\begin{proof}
En effet, en vertu de (\hyperref[0.13.5.6-fr]{13.5.6}), la condition
\emph{a)} équivaut à dire que le système projectif
$(\mathrm E_\infty^{pq}(A_k))_{k\in\mathbb Z}$ est essentiellement
constant pour $p+q=n$, et d'autre part
$\operatorname{gr}^p(\mathrm R^n\mathrm T(A_k))$ est canoniquement
isomorphe à $\mathrm E_\infty^{p,n-p}(A_k)$, donc il résulte de
(\hyperref[0.13.4.3-fr]{13.4.3}) que \emph{a)} et \emph{b)} sont
équivalentes~; l'isomorphisme
(\hyperref[0.13.7.4.1-fr]{13.7.4.1}) n'est autre que
(\hyperref[0.13.5.6.1-fr]{13.5.6.1}) appliqué au cas envisagé ici.
\end{proof}

\begin{corollary}[13.7.5]
\label{0.13.7.5-fr}
\emph{Soit $\mathcal F=(\mathcal F_k)_{k\in\mathbb N}$ un système
projectif de faisceaux de groupes abéliens satisfaisant aux conditions
\emph{(i)}, \emph{(ii)} et \emph{(iii)} de
(\hyperref[0.13.3.1-fr]{13.3.1}) et soit
$\mathcal F=\varprojlim_k\mathcal F_k$. Supposons que, pour le foncteur
$\mathcal G\mapsto\Gamma(X,\mathcal G)$, le système projectif
$(\mathrm E_r^{pq}(\mathcal F))_{r\in\mathbb Z}$ soit essentiellement
constant pour tout couple $(p,q)$ tel que $p+q=n$ ou $p+q=n+1$.
Considérons sur $H^{n+1}(X,\mathcal F)$ la filtration définie par
$\mathrm F^p(H^{n+1}(X,\mathcal F))=
\operatorname{Ker}(H^{n+1}(X,\mathcal F)\to
H^{n+1}(X,\mathcal F_{p-1}))$. On a alors un isomorphisme canonique}
\begin{equation}
\label{0.13.7.5.1-fr}
\tag{13.7.5.1}
\operatorname{gr}^p(H^{n+1}(X,\mathcal F))
\xrightarrow{\sim}\mathrm E_\infty^{p,n-p+1}(\mathcal F)
\qquad\emph{pour tout $p\in\mathbb Z$.}
\end{equation}
\end{corollary}
\begin{proof}
